% Source-derived chapter 03: Local Abel-Jacobi morphisms
% Original source: tates-thesis-de-rham-setting
\chapter{Local Abel-Jacobi morphisms}
\label{tates-thesis-de-rham-setting-chapter-03}
\source{tates-thesis-de-rham-setting}

\begin{remark}[Source section]
\label{tates-thesis-de-rham-setting-chapter-03-source}
\source{tates-thesis-de-rham-setting}
\source{tates-thesis-de-rham-setting}
This chapter reproduces the corresponding section of the retrieved
LaTeX source, including its definitions, statements, examples, and
proofs. It has not yet been translated into Lean.
\notready
\end{remark}

% The source section title was: Local Abel-Jacobi morphisms
\label{s:aj}
\source{tates-thesis-de-rham-setting}

\subsection{}

In this section, we collect some standard results
about local Abel-Jacobi maps and (simplified) 
Contou-Carr\`ere pairings for use in \S \ref{s:weyl}.
This material is standard and included for the reader's
convenience.

\subsection{}

Let $\widehat{\sD} \in \IndSch$ denote the formal disc
$\Spf(k[[t]]) \coloneqq \colim_n \Spec(k[[t]]/t^n)$.

Let $\widehat{\Delta}:\widehat{\sD} \to \widehat{\sD} \times \sD$ denote the
graph of the canonical map $\widehat{\sD} \to \sD$. 
Clearly this map is a closed embedding (in particular, schematic).

We define: 
%
\[
\sO(-\widehat{\Delta}) \coloneqq \Ker\big(
\sO_{\widehat{\sD} \times \sD} \to \widehat{\Delta}_*(\sO_{\widehat{\sD}})\big) 
\in \QCoh(\widehat{\sD} \times \sD).
\]

\noindent One readily checks 
that $\sO(-\widehat{\Delta})$ is a line bundle on 
$\widehat{\sD} \times \sD$: in fact, 
a choice of coordinate on the disc 
gives a trivialization of it. 
Moreover, as:
%
\[
\widehat{\sD}
\underset{\widehat{\sD} \times \sD }{\times}
(\widehat{\sD} \times \o{\sD} ) = \emptyset
\]

\noindent there is a canonical trivialization
of 
$\sO(-\widehat{\Delta})|_{\widehat{\sD} \times \o{\sD}}$. 

Thus, we obtain a map:
%
\[
\AJ^{-1}:\widehat{\sD} \to \Gr_{\bG_m}.
\] 

\noindent We 
define the \emph{local Abel-Jacobi map} 
$\AJ$ as $\AJ^{-1}$ composed
with the inversion map $\Gr_{\bG_m} \to \Gr_{\bG_m}$.

Explicitly, we have:
%
\[
\begin{gathered}
\AJ: \widehat{\sD} \to \Gr_{\bG_m} \\
\tau \mapsto (\sO_{\sD}(\tau),1)
\end{gathered}
\]

\noindent where $\sO_{\sD}(\tau)$ is the line bundle on 
the disc with sections having at worst simple
poles at $\tau \in \widehat{\sD} \subset \sD$
and the evident trivialization ``$1$" on $\o{\sD}$.

\subsection{}\label{ss:dlog-aj}
\source{tates-thesis-de-rham-setting}

We now consider a closely related map to the Abel-Jacobi map.

We have an evident evaluation map:
%
\begin{equation}\label{eq:ev-a1}
\source{tates-thesis-de-rham-setting}
\widehat{\sD} \times \fL^+ \bA^1 \to \bA^1
\end{equation}

\noindent that is linear in the second coordinate.
Using the non-degeneracy of the residue pairing, 
we deduce that there is a unique map:
%
\begin{equation}\label{eq:dlog-aj}
\source{tates-thesis-de-rham-setting}
\widehat{\sD} \to
\fL^{pol} \bA^1 dt.
\end{equation}

\noindent such that the composition:
%
\[
\widehat{\sD} \times \fL^+ \bA^1 
\xar{\eqref{eq:dlog-aj} \times \id}
\fL^{pol} \bA^1 dt \times 
\fL^+ \bA^1 \xar{(f,\omega) \mapsto \Res(f\omega)} 
\bA^1
\]

\noindent recovers the evaluation map above.

\begin{rem}
\source{tates-thesis-de-rham-setting}
\notready\label{r:dlog-aj-fmla}
\source{tates-thesis-de-rham-setting}

Explicitly, \eqref{eq:dlog-aj} is given in coordinates
by:
%
\[
(\tau \in \widehat{\sD}) \mapsto \sum_{i = 0}^{\infty}
\frac{\tau^i}{t^{i+1}} dt.
\]

\end{rem}

We have the following elementary calculation.

\begin{lem}
\source{tates-thesis-de-rham-setting}
\notready\label{l:dlog-aj}
\source{tates-thesis-de-rham-setting}

The map \eqref{eq:dlog-aj} coincides with the composition:
%
\[
\widehat{\sD} \xar{\AJ} \Gr_{\bG_m} 
\xar{-d\log} \fL^{pol} \bA^1 dt.
\]

\end{lem}

\begin{proof}

Choose coordinates as in Remark \ref{r:dlog-aj-fmla}.
The map $\AJ:\widehat{\sD} \to \Gr_{\bG_m}$ lifts to a 
map to $\fL\bG_m$ via $\tau \mapsto \frac{1}{t-\tau}$.
We then have:
%
\[
d\log \frac{1}{t-\tau} = 
\frac{d\frac{1}{t-\tau}}{\frac{1}{t-\tau}} = 
-\frac{dt}{t-\tau} = -\frac{1}{t}\sum_{i = 0}^{\infty} \frac{\tau^i}{t^i} dt
\]

\noindent Comparing to Remark \ref{r:dlog-aj-fmla}, we obtain the 
claim.

\end{proof}

\begin{rem}
\source{tates-thesis-de-rham-setting}
\notready

Following the lemma, we denote the above map by 
$-d\log \AJ$, and its additive inverse by $d\log \AJ$.

\end{rem}


\subsection{Compatibility with truncations}\label{ss:aj-trun}
\source{tates-thesis-de-rham-setting}

We use the following observation. We remind that
$\sD^{\leq n} \coloneqq \Spec(k[[t]]/t^n)$.

Clearly $d\log \AJ$ maps $\sD^{\leq n}$ maps 
into $\fL^{pol,\leq n} \bA^1 dt$. 
Therefore, by definition (see \S \ref{ss:loop-gr-trun}),
we deduce that $\AJ$ maps $\sD^{\leq n}$ into 
$\Gr_{\bG_m}^{\leq n}$.

\subsection{The positive Grassmannian}\label{ss:pos-gr}
\source{tates-thesis-de-rham-setting}

Next, we define: 
%
\[
\fL^{pos} \bG_m \coloneqq
\fL \bG_m \underset{\fL \bA^1} {\times} \fL^+ \bA^1.
\]

\noindent There is an evident commutative monoid
structure on $\fL^{pos} \bG_m$ and homomorphism
$\fL^+ \bG_m \to \fL^{pos} \bG_m$.

We then define:
%
\[
\Gr_{\bG_m}^{pos} \coloneqq \fL^{pos} \bG_m/\fL^+ \bG_m =
\Ker(\fL^+\bA^1/\bG_m \to \fL \bA^1/\bG_m).
\]

Explicitly, $\Gr_{\bG_m}^{pos}$ parametrizes
the data of $\sL$ a line bundle on the disc and
$\sigma \in \Gamma(\sD,\sL)$ a (regular) section
such that $\sigma|_{\o{\sD}}$ trivializes $\sL$.

\subsection{}

For $n \geq 0$, we define: 
%
\[
\Sym^n \widehat{\sD} \coloneqq \Spf(k[[t_1,\ldots,t_n]])^{S_n} 
\in \IndSch.
\]

\noindent for $S_n$ the symmetric group.
We then let: 
%
\[
\Sym\widehat{\sD} \coloneqq 
\coprod_n \Sym^n \widehat{\sD} \in \IndSch.
\]

\subsection{}

As for usual smooth curves, $\Sym\widehat{\sD}$ parametrizes
effective Cartier divisors on $\sD$ supported on $\widehat{\sD}$.
Hence, we obtain an isomorphism:
%
\[
\Sym\widehat{\sD} \isom \Gr_{\bG_m}^{pos}
\]

\noindent such that the composition:
%
\[
\widehat{\sD} = \Sym^1 \widehat{\sD} \to 
\Sym\widehat{\sD} \simeq \Gr_{\bG_m}^{pos}
\subset \Gr_{\bG_m}
\]

\noindent is $\AJ$.

This is an isomorphism of commutative monoids for the
evident product on $\Sym(\widehat{\sD})$.

\subsection{Log Contou-Carr\'ere}

\subsubsection{}

We now observe the following.

There is a unique pairing:
%
\[
\pair{-}{-}:\Gr_{\bG_m}^{pos} \times \fL^+ \bA^1 \to \bA^1
\]

\noindent that is bilinear for multiplicative monoid
structures on both sides and whose restriction 
to $\widehat{\sD} \times \fL^+ \bA^1$ is 
\eqref{eq:ev-a1}. 

Explicitly, for $D \in \Gr_{\bG_m}^{pos}$ an effective 
Cartier divisor
on $\sD$ supported on $\widehat{\sD}$ and 
$f\in \fL^+ \bA^1$, we have:
%
\[
\pair{D}{f} = f(D) \coloneqq \Nm_D(f).
\]

\subsubsection{}

There are some variants of the pairing above.

First, note that 
$\pair{-}{-}|_{\Gr_{\bG_m}^{pos} \times \fL^+ \bG_m}$
maps into $\bG_m$. As $\Gr_{\bG_m}$ is the group completion
of $\Gr_{\bG_m}^{pos}$ (in the evident sense), we obtain a 
pairing:
%
\[
\pair{-}{-}:\Gr_{\bG_m} \times \fL^+ \bG_m \to 
\bG_m.
\]

\noindent This pairing is compatible with Contou-Carr\'ere's
pairing \cite{cc} in the evident sense.

\subsubsection{}

Next, let $\Sym^m \sD^{\leq n} \coloneqq
\Spec((k[[t]]/t^n)^{\otimes m,S_m})$, which is an
affine scheme. We let $\Gr_{\bG_m}^{pos,\leq n}\coloneqq 
\coprod_m \Sym^m \sD^{\leq n}$.

The evident map:
%
\[
\Gr_{\bG_m}^{pos,\leq n} \to \coprod_m \Sym^m \widehat{\sD} \to \Gr_{\bG_m}
\]

\noindent maps into $\Gr_{\bG_m}^{\leq n}$. Indeed,
this map is a map of monoids, so the same is
true of its composition with $d\log:\Gr_{\bG_m} \to 
\fL^{pol} \bA^1 dt$. As $\sD^{\leq n}$ maps into
$\fL^{pol,\leq n} \bA^1 dt$ (cf. \S \ref{ss:aj-trun}),
we obtain the claim by definition of $\Gr_{\bG_m}^{\leq n}$.


For similar reasons, the pairing:
%
\[
\pair{-}{-}:
\Gr_{\bG_m}^{pos,\leq n} \times \fL^+ \bA^1 \to \bA^1
\]

\noindent factors through a pairing:
%
\[
\Gr_{\bG_m}^{pos,\leq n} \times \fL_n^+ \bA^1 \to \bA^1 
\]

\noindent that we also denote by $\pair{-}{-}$.

\begin{rem}
\source{tates-thesis-de-rham-setting}
\notready\label{r:gr-n-pos-gp-cpl}
\source{tates-thesis-de-rham-setting}

The base-point $0 \in \widehat{\sD}$ gives a point of 
$\Gr_{\bG_m}^{pos,\leq n}$ (of degree 1), hence 
a translation map $T:\Gr_{\bG_m}^{pos,\leq n} \to \Gr_{\bG_m}^{pos,\leq n}$.
There is a natural map:
%
\[
\colim \, (\Gr_{\bG_m}^{pos,\leq n} \xar{T} \Gr_{\bG_m}^{pos,\leq n} \xar{T} \ldots)
\to \Gr_{\bG_m}^{\leq n}
\]

\noindent that we claim is an isomorphism. 
Indeed, the assertion is evident at the level of $k$-points (both sides 
are $\bZ$), so it suffices to check the assertion on tangent complexes,
where it is straightforward (both sides naturally identify with 
$t^{-n+1}k[[t]]/k[[t]]$). 

\end{rem}


\subsubsection{}\label{ss:pairing-map}
\source{tates-thesis-de-rham-setting}

We will use the following constructions (in a quite simple setting).

Let $V$ be a vector space. We consider $V$ as a prestack\footnote{If $Y$ is
coconnective, this prestack is a (DG) indscheme.}
with functor of points $\Spec(A) \mapsto \Omega^{\infty}(A\otimes V) \in \Gpd$.

Now suppose $X$ is an ind-proper indscheme equipped with a morphism
$\vph:X \to V$. We obtain a canonical ``integration" map:
%
\begin{equation}\label{eq:integration}
\source{tates-thesis-de-rham-setting}
\Gamma^{\IndCoh}(X,\omega_X) \to V.
\end{equation}

\noindent Indeed, our map $\vph$ is (tautologically) equivalent to
a morphism $\sO_X \to V\otimes \sO_X \in \QCoh(X)$. Tensoring with the dualizing
sheaf on $X$, we obtain a morphism $\omega_X \to V \otimes \omega_X \in \IndCoh(X)$,
or what is the same, a map $p^!(k) \to p^!(V)$ for $p:X \to \Spec(k)$ the projection.
By ind-properness and adjunction, 
we obtain a map $p_*^{\IndCoh}p^!(k) \to V$ as desired. 

\begin{rem}
\source{tates-thesis-de-rham-setting}
\notready\label{r:upsilon-ff}
\source{tates-thesis-de-rham-setting}

Using \cite{grbook} \S II.3 Lemma 3.3.7, one finds that this is a bijection;
i.e., specifying a map \eqref{eq:integration} is equivalent to giving
$X \to V$.

\end{rem}

We also have the following variant.
Suppose we are also given $Y$ a prestack and
$\pi:X \times Y \to \bA^1$ a map.

The map $\pi$ can be encoded by a morphism $X \to \Gamma(Y,\sO_Y)$ (thinking of the target
as a prestack as above). Therefore, 
we obtain a comparison map:
%
\[
c_{\pi}:\Gamma^{\IndCoh}(X,\omega_X) \to \Gamma(Y,\sO_Y) \in 
\Vect
\]

\noindent by the above. (As in Remark \ref{r:upsilon-ff}, 
one can actually recover $\pi$ from $c_{\pi}$.) 

\subsubsection{}

We now have the following non-degeneracy assertion for our
pairings.

\begin{prop}
\source{tates-thesis-de-rham-setting}
\notready\label{p:cc-log}
\source{tates-thesis-de-rham-setting}

Each of the morphisms:
%
\[
\begin{gathered}
\Gamma^{\IndCoh}(\Gr_{\bG_m}^{pos},\omega_{\Gr_{\bG_m}^{pos}}) 
\to \Gamma(\fL^+ \bA^1,\sO_{\fL^+ \bA^1}) \\
\Gamma^{\IndCoh}(\Gr_{\bG_m},\omega_{\Gr_{\bG_m}}) 
\to \Gamma(\fL^+ \bG_m,\sO_{\fL^+ \bG_m}) \\
\Gamma^{\IndCoh}(\Gr_{\bG_m}^{pos,\leq n},\omega_{\Gr_{\bG_m}^{pos,\leq n}}) 
\to \Gamma(\fL_n^+ \bA^1,\sO_{\fL_n^+ \bA^1}) \\
\Gamma^{\IndCoh}(\Gr_{\bG_m}^{\leq n},\omega_{\Gr_{\bG_m}^{\leq n}}) 
\to \Gamma(\fL_n^+ \bG_m,\sO_{\fL_n^+ \bA^1})
\end{gathered}
\]

\noindent coming from our pairings and \S \ref{ss:pairing-map}
is an isomorphism.

Moreover, each isomorphism is a map of commutative algebras,
where the left hand sides are commutative algebras using
the commutative monoid structures on the relevant spaces,
while the right hand sides are commutative algebras 
as functions on schemes.

\end{prop}

\begin{proof}

Each of these maps is a map of commutative algebras
because of the bilinearity of the pairings.

Next, we have:
%
\[
\begin{gathered}
\Gamma^{\IndCoh}(\Gr_{\bG_m}^{pos,\leq n},\omega_{\Gr_{\bG_m}^{pos,\leq n}}) =
\underset{m}{\oplus} \,
\Gamma^{\IndCoh}(\Sym^m\sD^{\leq n},\omega_{\Sym^m\sD^{\leq n}})
= \\
\underset{m}{\oplus} \, ((k[[t]]/t^n)^{\otimes m,S_m})^{\vee}
= \Sym((k[[t]]/t^n)^{\vee})
\end{gathered}
\]

\noindent where the second equality is by finiteness of
$\Sym^m \sD^{\leq n}$.

We similarly have 
$\fL_n^+ \bA^1 = \Spec(\Sym((k[[t]]/t^n)^{\vee}))$ by 
definition. Under these identifications, it suffices to
show that the map:
%
\[
\Sym((k[[t]]/t^n)^{\vee}) = 
\Gamma^{\IndCoh}(\Gr_{\bG_m}^{pos,\leq n},\omega_{\Gr_{\bG_m}^{pos,\leq n}}) 
\to \Gamma(\fL_n^+ \bA^1,\sO_{\fL_n^+ \bA^1}) =
\Sym((k[[t]]/t^n)^{\vee})
\]

\noindent from the proposition is the identity: 
it suffices to check this on
the generators of the left hand side (as we have a map
of commutative algebras), and there it follows 
by construction of the pairing.

Passing to the colimit as $n \to \infty$, we recover
the first map under consideration.

The second and fourth cases follow from Remark \ref{r:gr-n-pos-gp-cpl}.

\end{proof}

\begin{rem}
\source{tates-thesis-de-rham-setting}
\notready

The assertion $\Gamma^{\IndCoh}(\Gr_{\bG_m},\omega_{\Gr_{\bG_m}}) 
\simeq \Gamma(\fL^+ \bG_m,\sO_{\fL^+ \bG_m})$ via this
pairing is standard from Contou-Carr\`ere \cite{cc}. 
The other assertions
are apparently less well-known,\footnote{For example,
they concern Cartier duality for commutative monoids,
which is less often considered than the group case.}
though they easy results.

\end{rem}

%\include{sect3}
